\section{Prepared experiments and stopping resources}
\label{sec:experiment}
All state, action and report spaces are standard Borel. An unknown index $\theta\in\Theta$ specifies a prepared law and positive normalized causal kernels (positivity means nonnegative measures, not a strictly positive density at every report). It is fixed for an entire execution; it may index an arbitrary history-dependent physical evolution. Policies see only past reports, legal physical commands and their own allowed coins. They never see $\theta$. All assertions below are for the declared experiment and test, not for a formal parameter manifold.

A call has a nonnegative cost bounded below by $c_*>0$. A predictable flag declares whether the call is productive. Commitment is allowed only after a productive report. Conditional on commitment, a fixed bounded future audit $Z\in H$ is executed, where $H\subset\R^k$ is compact and convex. This audit is not observed before the prediction is committed. The loss is $\|Z-r\|^2$, $r\in H$. Each productive call has a potential audit specified by the same kernel, whether or not the observer commits there. A fresh stopping coin is independent of that audit conditional on the current physical report and history. All other informative components of the current report must be included in the conditional laws below.

The cost $N$ in the core theorem counts preparation, failed gate calls, productive calls, and all calls preceding commitment. A separately required final physical audit has its own explicit cost. Thus a protocol with one initial preparation call and one final audit uses total raw budget $N+2$, not $N$. Time, bit operations and scratch space are not identified with raw-call count.

\subsection*{The stopping interface and retained labels}
A checkpoint procedure may retain the entire precommit history, but its committed value belongs to one fixed dictionary of at most $m$ elements. This convention counts only its terminal dictionary, not its full-history working memory. Its physical acquisition and audit costs are still charged.

A finite-state observer has an alphabet $S$ of at most $M$ labels and one fixed parameter-independent program. At every intercall boundary its only observer memory is the retained label. Its next physical command distribution is a fixed function of that label. In particular, the physical port identity, elapsed time, prior report and subprogram index are not additional free inputs. During a call the current command, its report and fresh independent coins are available; the post-report transition can depend on them. No such coin, report, real register or clock persists unless encoded in $S$. A fixed terminal decoder $r:S\to H$ reads only the committed label; it never reads the report or the audit. The two stopping interfaces are as follows.

\begin{definition}[Event and state stopping]\label{def:stopping}
In the \emph{event-stopped} interface, a post-report kernel emits $(s',e)\in S\times\{0,1\}$. The value $e=1$ is permitted only after a productive call and commits exactly once: the experiment stops, the label $s'$ is retained, and the decoder returns $r(s')$. No next command is executed. The commitment event belongs to the public stopped transcript; it is not an extra stored bit available during subsequent computation. On branches with $e=0$, execution continues at $s'$. A label may therefore be used both live and terminal on different branches.

In the \emph{state-stopped} interface there is instead a fixed map $h_*:S\to\{0,1\}$. The initial label has type zero. A transition commits if and only if $h_*(s')=1$, and entry into type one is permitted only after a productive report. Type-one labels are globally terminal and absorbing: they have no executed command row. Type-zero labels are live. The halt status is thus completely determined by the retained state, and live and terminal labels are disjoint.
\end{definition}

Let $L_\gamma$ be the set of labels reached before physical calls and $T_\gamma$ the set reached on commitment, under the relevant family of path laws. The program's role counts are $\ell_\gamma=|L_\gamma|$ and $t_\gamma=|T_\gamma|$, with overlap $o_\gamma=|L_\gamma\cap T_\gamma|$. Its actual alphabet count is $\ell_\gamma+t_\gamma-o_\gamma$, whereas its number of distinct decoder values is at most $t_\gamma$. A mandatory lower on $\ell_\gamma$ does not imply a subtraction from $M$ unless overlap is excluded. This is the distinction used below, not an unmentioned memory convention.

Under robust almost-sure commitment and robust expected core cost at most $N$, write the infimal worst-case risks as $R_\cp^\Theta(N,m)$, $R_\evt^\Theta(N,M)$ and $R_\st^\Theta(N,M)$. An empty legal class has value $+\infty$. Ideal Borel programs may have fixed real constants; finite descriptions, evaluation time and scratch space are treated separately. A randomly selected program used for the entire execution is not a free persistent seed. Both stopping interfaces are legitimate resources, but neither is silently identified with the other.
